\subsubsection{Systems}\label{subsec:systems-riemann-problems}

In this section we state a collection of results that allow us to solve the Riemann problem for systems of equations, but some of our proofs are only for the special case of the $p$ system. This allows us to keep our treatment fairly brief and concrete while displaying most of the ideas involved in the more general proofs.

For the single conservation law we were able to connect any pair of left and right states using a single wave, either a shock or a rarefaction wave. In higher dimensions, we will have to use intermediate states and several different waves to make the connection. However, as a first step, we will see what left and right states can be "hooked up" using a single shock or rarefaction wave.

\paragraph*{Shock waves}

We begin by considering the possibility of using a single shock wave to connect the left and right states. Thus, we have to ask the question: given $\bm{u}^l$, what states $\bm{u}^r$ satisfy the Rankine-Hugoniot condition \eqref{eq:rankine-hugoniot-jump-condition} and the Lax shock condition \eqref{eq:lax-shock-condition}? The answer is that, emanating from each point $\bm{u}^l$ in state space, there are $n$ \emph{shock curves} that describe the possible right states that can be connected by a single shock. More specifically, we have the following theorem.

\begin{Theorem}
\label{thm:shock-curves-existence}
Suppose that \eqref{eq:conservation-law-system} is a strictly hyperbolic system of conservation laws defined on a region $\Omega\subset\mathbb{R}^n$ of state space. Then for any $\bm{u}^l\in\Omega$ there exist $n$ open intervals $I_k$ containing $0$ and $n$ one-parameter families of states $\hat{\bm{u}}_k(\epsilon)$ and shock speeds $\hat s_k(\epsilon)$ defined on $\epsilon\in I_k$ such that
\begin{equation}\label{eq:shock-curve-base-state}
\bm{u}_k(0)=\bm{u}^l
\end{equation}
and such that for $\epsilon\in I_k$, $\hat{\bm{u}}_k(\epsilon)$ and $\hat s(\epsilon)$ satisfy the Rankine-Hugoniot condition
\[
\hat s(\epsilon)[\bm{u}^l-\hat{\bm{u}}_k(\epsilon)]=\bm{f}(\bm{u}^l)-\bm{f}(\hat{\bm{u}}_k(\epsilon)).
\]
Furthermore, if the $k$th characteristic field is genuinely nonlinear, then the parameterization can be chosen so that
\begin{equation}\label{eq:shock-curve-initial-tangent}
\hat{\bm{u}}_k'(0)=\bm{r}_k(\bm{u}^l),
\end{equation}

\begin{equation}\label{eq:shock-curve-initial-speed}
\hat s(0)=\lambda_k(\bm{u}^l),
\end{equation}

\begin{equation}\label{eq:shock-curve-initial-speed-derivative}
\hat s'(0)=1/2,
\end{equation}
where the prime $'$ refers to differentiation with respect to $\epsilon$. Moreover, with this parameterization, the Lax shock conditions hold if and only if $\epsilon<0$.
\end{Theorem}

We will not prove this theorem in general, but will instead calculate the shock curves explicitly for the $p$ system. To ensure strict hyperbolicity and genuine nonlinearity we will assume $p'<0$ and $p''>0$. We can also either assume that $p$ is defined on all of $\mathbb{R}$ or make appropriate restrictions on the states chosen below. Thus, we take any admissible $\bm{u}^l:=(w^l,v^l)^t$ and suppose $\hat{\bm{u}}:=(\hat w,\hat v)^t$ is connected to $\bm{u}^l$ by one of the two shock curves whose existence was asserted in the theorem. In this case the Rankine-Hugoniot condition reduces to
\begin{equation}\label{eq:p-system-shock-volume-rh}
s(w^l-\hat w)=-(v^l-\hat v),
\end{equation}

\begin{equation}\label{eq:p-system-shock-velocity-rh}
s(v^l-\hat v)=p(w^l)-p(\hat w).
\end{equation}
By eliminating $s$ from these equations we get
\begin{equation}\label{eq:p-system-shock-eliminated-speed}
(v^l-\hat v)^2=(\hat w-w^l)(p(w^l)-p(\hat w)).
\end{equation}


Since $p'<0$, there are two curves of solutions, defined for all $\epsilon$ in the domain of $p$.
\begin{equation}\label{eq:p-system-first-shock-curve}
\begin{aligned}
S_1:\quad\hat{\bm{u}}_1(\epsilon)&=\begin{pmatrix}\hat w_1(\epsilon)\\\hat v_1(\epsilon)\end{pmatrix}\\
&:=\begin{pmatrix}\epsilon\\v^l+\operatorname{sgn}(\epsilon-w^l)\sqrt{(\epsilon-w^l)(p(w^l)-p(\epsilon))}\end{pmatrix},
\end{aligned}
\end{equation}

\begin{equation}\label{eq:p-system-second-shock-curve}
\begin{aligned}
S_2:\quad\hat{\bm{u}}_2(\epsilon)&=\begin{pmatrix}\hat w_2(\epsilon)\\\hat v_2(\epsilon)\end{pmatrix}\\
&:=\begin{pmatrix}\epsilon\\v^l-\operatorname{sgn}(\epsilon-w^l)\sqrt{(\epsilon-w^l)(p(w^l)-p(\epsilon))}\end{pmatrix}.
\end{aligned}
\end{equation}
The corresponding shock speeds are
\begin{equation}\label{eq:p-system-first-shock-curve-speed}
\hat s_1(\epsilon):=-\sqrt{\dfrac{p(w^l)-p(\epsilon)}{\epsilon-w^l}},
\end{equation}

\begin{equation}\label{eq:p-system-second-shock-curve-speed}
\hat s_2(\epsilon):=\sqrt{\dfrac{p(w^l)-p(\epsilon)}{\epsilon-w^l}}.
\end{equation}


\begin{figure}[H]
\centering
\includegraphics[width=.78\linewidth]{figures/p-system-shock-curves-source.png}
\caption{}\label{fig:p-system-shock-curves}
\end{figure}

Only half of each curve satisfies the Lax shock conditions. Conditions \eqref{eq:p-system-back-shock-volume-jump} and \eqref{eq:p-system-front-shock-volume-jump} imply that any right state of a 1-shock would have to lie on the curve
\begin{equation}\label{eq:p-system-admissible-first-shock-states}
\bm{u}^r=\hat{\bm{u}}_1(\epsilon),\quad\epsilon<w^l,
\end{equation}
and any right state of a 2-shock would have to lie on the curve
\begin{equation}\label{eq:p-system-admissible-second-shock-states}
\bm{u}^r=\hat{\bm{u}}_2(\epsilon),\quad w^l<\epsilon.
\end{equation}
The reader is asked to verify that these curves can be reparameterized so that they satisfy the stated initial conditions; and more importantly, that these states and the corresponding shock speeds satisfy \eqref{eq:p-system-back-shock-speeds} and \eqref{eq:p-system-front-shock-speeds} (cf. Problem~\ref{ex:p-system-shock-reparametrization}).

Pictorially, we see that emanating from each left state $\bm{u}^l$ we have the two shock curves $S_1$ and $S_2$. Shocks with negative speed (sometimes called slow shocks or back-shocks) lie along $S_1$; shocks with positive speed (fast shocks or front-shocks) lie along $S_2$.
